\section{Native rank-two support on the actual selected host}
\label{sec:cubic-rank-two-native-support}

This section supplies the rank-two input retained in AC-G106-even.
The host is the actual selected host, which need not be abelian
over \(\mathbb Q\). We construct the native \path{Setup}
directly. In particular the proof does not apply the packaged
\path{CRealEx} theorem after omitting its \path{GalLink.ab} field.

\paragraph{Retained arithmetic instance.}
Keep one completed AC-G78--AC-G93-fill instance at \(p=269\),
and then any sufficiently high finite cyclotomic layer of its
CM host \(H/F\), \(F=H^+\). In particular keep the actual
old normalized endpoint representation \(V\) supplied by
AC-G68-point and the good-twist construction: it is continuous,
integral over one finite coefficient field, has determinant
\(\eta_K\), and has an ordinary cuspidal modular realization
of weight \(2e+1\), with the exact original finite characters
and bad level retained. Its restriction to \(G_K\) is
irreducible. The no-quadratic-self-twist argument of
AC-K29-nonCM is verified again below at the present prime,
so no excluded-prime scope is imported from an earlier statement.
The concrete progression
\[
                         e\equiv267\pmod {268}
\tag{G106N.1}
\]
is sufficient for the completed host instance used here.
No existence of such an old point is inferred solely from a
residual tuple. This is the same starting arithmetic instance
as AC-G78 and AC-G87.

Retain the published cyclic base-change and automorphic
attachment inputs used in that instance, the exact native
\path{LevEx}, \path{ProModularRingExists}, \path{LamEx},
and \path{ModU}, and local and global reciprocity. The old
automorphy input and its original coefficient hypotheses remain
those of AC-G68-point. The finite coefficient ring
\(\mathcal O\) in the present native construction is the
already chosen large coefficient ring of the host instance;
it contains the values of the character \(\kappa_2\) of
AC-G106-twist. Enlarging it once, before making the present
rings, is allowed. This does not change or remove the original
\(e_O=1\) premise used to obtain \(V\).

Use the rational-character frame \(\mathcal F_\gamma\)
and the rings \(B_2\subset A_2\subset R_2\) of
AC-G106-frame and AC-G106-corner. Thus
\(\gamma=\bar\psi\bar\epsilon^e\) is the actual rational
residual character, the residual ordered tuple is
\((\gamma,\gamma^{-1})\), and the weights are \((1,-1)\).
The symbol \(\epsilon\) in this section is the native
cyclotomic character: on geometric Frobenius at a place of
norm \(q\), it has value \(q^{-1}\). This is the convention
in the literal native \path{GFrame.eps} and
\path{GFrame.starOp}.

\begin{theorem}[AC-G106-native: actual native rank-two support and uniform even support]
\label{AC-G106-native}
Under the retained inputs above, at every sufficiently high
finite cyclotomic layer there is an actual native rank-two
ordinary determinant ring \(\mathscr B_2\), an actual native
ideal \(J_2\subset\mathscr B_2\), and a surjection
\[
             \Phi_2:\mathscr B_2\longtwoheadrightarrow B_2
\tag{G106N.2}
\]
such that, with \(M_2=\Phi_2(J_2)B_2\),
\[
 \begin{gathered}
 M_2\subsetneq B_2,\qquad
 \dim_k B_2/(M_2+\lambda_2)<\infty,\qquad
 \sqrt{M_2+\lambda_2}=\mathfrak m_{B_2},\\
 M_2\subset P\cap B_2
 \quad\text{for every prime }P\subset A_2
                         \text{ off both triangular loci}.
 \end{gathered}
\tag{G106N.3}
\]
Here \(\lambda_2\) is the literal rational-frame weight-fiber
ideal, generated by the uniformizer and
\(\lambda_{v,i}(g)-1\) for
\(g\in I_v\cap\ker\gamma\). In particular this is
the exact source \path{OffFin}, with one ideal for all such
primes, not an assertion only on the component through the
automorphic seed.

Consequently the ideal \(M_E\subset B_\theta\) constructed
in AC-G106-even satisfies the actual changed-tuple
\path{FinLam} condition and is contained in every prime of
\(A_\theta\) which is both even and off triangular.
This supplies the uniform even-support input used in the
normalized propagation argument. All native arithmetic
inputs displayed above remain explicit; no uniform numerical
bound on a host, a class group, or a Selmer length is asserted.
\end{theorem}

\begin{proof}
\proofstep{1. The native group and the genuine rank-two twist link.}
Use the actual \path{HidaSetting} \(s\), its full local
groups, and the \path{Lev} with \(V.S'=s.S\) from
AC-G97. Write
\[
  \mathcal G_S=G_F/N_S,\qquad
  \mathcal G_{S,0}=\ker\eta_S,\qquad G_\gamma=G_H/N_\gamma.
\tag{G106N.4}
\]
The subgroup \(N_\gamma\) is the closed normal subgroup
generated by \(I_w\cap\ker\gamma\), \(w\nmid p\),
including all conjugate places. Every inertia group outside
\(S\) belongs to this family. Therefore \(N_S\subset
N_\gamma\), and the quotient map gives a continuous
surjection \(\pi:\mathcal G_{S,0}\to G_\gamma\).
It commutes with complex conjugation, the cyclotomic
character, and all chosen local maps. On coefficient-prime
local groups the comparison is the identity between the
same full decomposition groups; it is not a comparison of
their images in the global quotient.

The character \(\kappa_2:G_H\to\mathcal O^\times\)
of AC-G106-twist is unramified outside \(3p\).
These places belong to \(S\), so it kills \(N_S\) and
descends to \(\mathcal G_{S,0}\). It need not descend
to \(G_\gamma\). The proved norm relation is exactly
\[
             \kappa_2(g)\kappa_2(cgc^{-1})
                          =\epsilon(g)^{-1}.
\tag{G106N.5}
\]
For \(i=0,1\), put
\[
  \bar t_i=(\gamma^{w_i}\circ\pi)\bar\kappa_2,
  \qquad (w_0,w_1)=(1,-1).
\tag{G106N.6}
\]
These are distinct finite residual characters with open
kernels. Their native polarization follows from
(G106N.5) and \(\gamma^c=\gamma\). More explicitly,
the native rank-two star sends \([g]\) to
\(\epsilon(g)^{-1}[(cgc^{-1})^{-1}]\); on the diagonal
representation (G106N.6), it is transpose conjugated by
\(J_2=\left(\begin{smallmatrix}0&1\\1&0\end{smallmatrix}\right)\).
This proves the polynomial-law polarization after every
coefficient base change. Take the local residual characters
to be its restrictions, and \path{th}=1.
By AC-G97, the literal native \path{mu} is the closed
local commutator. All these characters kill it. This
checks every field of the native rank-two \path{Resid}.

We check the remaining \path{TwLink} fields explicitly.
For away inertia, \(\pi(I_w)\) is a finite group of order
prime to \(p\), since it factors through \(\gamma(I_w)\).
The inertia image of \(\kappa_2\) has exponent dividing
six. A common multiple of these two prime-to-\(p\)
orders gives the exact \path{tame} field. At the auxiliary
place \(v_0\), \(\gamma\) is unramified and
\(\kappa_2\) is unramified, so both are trivial on
the native \path{I0}. This is \path{ramκ}.
The source \path{ker} field says that any closed normal,
conjugation-stable subgroup containing the away inertia
elements with \(\gamma=1\) and \path{I0} contains
\(\ker\pi\). This is precisely the definition of
\(N_\gamma/N_S\); no larger inertia kernel has been
silently identified with a smaller one.
The local isomorphisms preserve inertia, and all continuity
and \path{chib}/\path{psib} identities are (G106N.4)--(G106N.6).
Thus this is the literal original rank-two \path{TwLink},
not the changed-tuple replacement used in rank three.

\proofstep{2. The native Pair, including the potentially larger native inertia kernel.}
For the universal rational rank-two representation \(r_2\)
and its local characters \(\lambda_{v,i}\), set
\[
 \widetilde r_2(g)=\kappa_2(g)r_2(\pi g),\qquad
 \widetilde\lambda_{v,i}(h)
       =\kappa_2(\phi_v h)\lambda_{v,i}(h).
\tag{G106N.7}
\]
The determinant law of this representation has coefficients
in \(B_2\). Indeed a coefficient of the determinant of
a finite sum of its matrices is a sum of zero-weight
entry monomials; scalar \(\kappa_2\)-values lie in
\(\mathcal O^\times\). The same coefficientwise
argument works over arbitrary coefficient algebras and
therefore defines the required degree-two polynomial law.
It is continuous in the topology of \(B_2\). Here is the
required topology check. In a finite homogeneous presentation
of the rank-two universal ring, a weight-zero monomial decomposes
into weight-zero factors and pairs of weights \(2,-2\).
Grouping these factors in the formal presentation, and then
passing to the quotient as in AC-G95, gives, for \(m\ge1\),
\[
       \mathfrak m_{R_2}^{\,3m+2}\cap B_2
                              \subset\mathfrak m_{B_2}^{\,m}.
\]
The exponent has harmless slack: each nontrivial atom here
has length at most two. The grading gives
\(\mathfrak m_{B_2}^mA_2\cap(A_2)_0=\mathfrak m_{B_2}^m\).
Consequently every continuous universal coefficient with
values in \(B_2\) is continuous in its actual adic topology.
This is the rank-two instance of the full coefficientwise
argument of AC-G97-pair. The reduction, local-character fields and
their \path{mu} conditions are exactly (G106N.6).

Polarity and (G106N.5) make the native star transpose
conjugation by \(J_2\), proving \path{Pair.P}.
For the ordinary relation the two factors are respectively
\(\kappa_2(h_i)(r_2(h_i)-\lambda_i(h_i)I)\).
Their ordered product is zero, by the defining rank-two
ordinary relation. The same holds after multiplying on
the left by any group matrix, giving \path{Pair.Ord}.

The literal native kernel is
\[
                    \operatorname{Ip}(w)
              =I_w\cap\ker\bar t_0\cap\ker\bar t_1.
\tag{G106N.8}
\]
It is not asserted to equal \(I_w\cap\ker\gamma\).
If \(g\) belongs to (G106N.8), then
\(\gamma(g)^2=1\) and
\(\bar\kappa_2(g)=\gamma(g)^{-1}=\gamma(g)\).
The matrix \(r_2(\pi g)\) has finite order prime to
\(p\), and its residue is the scalar \(\gamma(g)I\).
The finite-order scalar lemma of AC-G106-corner gives
\(r_2(\pi g)=[\gamma(g)]I\).
The finite prime-to-\(p\) character value
\(\kappa_2(g)\) is its own Teichmuller lift, so their
product is one. Thus \(\widetilde r_2(g)=I\).
This proves \path{Pair.I} over the full possibly
nonreduced universal ring. In particular the case
\(\gamma(g)=\bar\kappa_2(g)=-1\) is included.

The ring \(B_2\) is an actual \path{IsCO} test ring.
The retained \path{ProModularRingExists} and the native
universal property now give \(\mathscr B_2\) and
(G106N.2) as a map, with the exact native \path{IsPhi}
identity. All the data of the source \path{Setup}
are supplied: AC-G106-frame provides \path{Ok} and
\path{Big} at sufficiently high layers, \(p\ne0\)
in \(\mathcal O\), and the rational \path{Space}
is the one just used.

\proofstep{3. Surjectivity on the actual invariant ring.}
Put \(C=\operatorname{im}\Phi_2\subset B_2\).
The native ring is complete Noetherian local with finite
residue field, hence compact. Its map is continuous and
local. Thus \(C\) is a closed complete local subring
of \(B_2\), with the same residue field \(k\), by the
compact quotient argument of AC-G102-surj.
Untwisting traces and local characters in (G106N.7) puts
every \(\operatorname{tr}r_2(g)\) and
\(\lambda_{v,i}(h)\) in \(C\).

Let \(Q=r_2(g_0)-r_2(g_0^\star)\).
Rank-two polarity gives directly
\[
 Q=\operatorname{diag}(a,-a),\qquad
 \bar a=\gamma(g_0)-\gamma(g_0)^{-1}\ne0.
\tag{G106N.9}
\]
The trace expansion puts
\(a^2=\tfrac12\operatorname{tr}Q^2\) in \(C\).
The simple residual root \(\bar a\) lifts by complete
local Hensel lifting to a root \(a_C\in C\).
In the ambient ring
\((a-a_C)(a+a_C)=0\), and \(a+a_C\) is a unit,
so \(a=a_C\). Hence \(a,a^{-1}\in C\), also in
the presence of nilpotents or coefficient-prime torsion.
The coordinate projectors are
\[
                   E_0=\frac{Q+aI}{2a},\qquad
                   E_1=\frac{aI-Q}{2a}.
\tag{G106N.10}
\]
It follows by expanding these expressions into group
matrices that
\[
 r_{ii}(g)=\operatorname{tr}(E_i r_2(g))\in C,
 \qquad
 r_{01}(g)r_{10}(h)
       =\operatorname{tr}(E_0r_2(g)E_1r_2(h))\in C.
\tag{G106N.11}
\]
Every zero-weight entry monomial is a product of the
diagonal entries and these upper--lower pairs: its upper
and lower factors have equal multiplicity and their
values commute. These monomials and the local characters
topologically generate \(B_2\). The image \(C\) is
closed, so \(C=B_2\). This proves the asserted
surjectivity without assuming reducedness.


\proofstep{4. The actual automorphic seed and its native arithmetic point.}
First \(V\) has no nontrivial quadratic self-twist on
\(G_{\mathbb Q}\). If \(V\simeq V\delta\) for a nontrivial
quadratic character, reduction of the trace identity for
\(\bar V^{\rm ss}=\eta_K\gamma\oplus\gamma^{-1}\)
gives \(\bar\delta=\eta_K\gamma^2\).
For clarity, on the nontrivial \(\delta\)-coset the sum
of the two residual characters is zero, so their ratio
is \(-1\); dividing two such evaluations makes their
ratio one on the kernel coset. This is exactly the
displayed character identity. Since \(\eta_K\) is
unramified at \(p\) and \(\gamma^2|_{I_p}\ne1\),
\(\delta|_{I_p}\ne1\). Apply the same two-coset argument
to the two characteristic-zero characters of the actual
ordinary flag on \(G_{K_p}\). It would make their ratio equal
to \(\delta|_{G_{K_p}}\). Their inertia ratio is
\(\epsilon^{2e}\), since \(\psi_p^2=1\), and has
infinite image. This contradicts the order two of
\(\delta\). The argument is valid at \(269\) itself.

The restriction to the actual \(G_H\) is now absolutely
irreducible. The initial field \(H_{\rm ab}\) in the
actual host construction is abelian over \(\mathbb Q\),
and contains \(K\). For an irreducible representation
and a normal subgroup of finite index, restriction is
semisimple: translates of a simple submodule sum to
the whole representation. Frobenius reciprocity for
the abelian quotient decomposes its restricted
endomorphism space into character-self-twist spaces.
Any self-twist of a rank-two representation has square
one by determinants. The preceding paragraph excludes
every nontrivial such character over \(\mathbb Q\).
Thus \(V|_{G_{H_{\rm ab}}}\) is absolutely irreducible.
At each subsequent cyclic step of degree \(p\), the same
argument excludes a nontrivial quotient self-twist,
because a character of order \(p\) cannot have square one.
These are exactly the abelian-then-cyclic-\(p\) extensions
in AC-G93-fill and its finite cyclotomic layers.
The argument applies after algebraic coefficient extension.

The retained cyclic automorphic base change gives a
cuspidal representation throughout this tower: its
noncuspidal alternative would be precisely such a
character self-twist. The retained unramified attachment
and Chebotarev identify that alternative with a self-twist
of the attached representation, so the preceding exclusion
applies at every step. Let \(\pi_0\) be the unitary
normalization on \(\operatorname{GL}_2(\mathbb A_H)\)
attached to \(V|_{G_H}\). Let \(Y\) be the algebraic
Hecke character of AC-G106-twist, restricted by the norm
from \(H\) to \(K\). Its infinity type is \((-1,0)\)
at the two distinguished embeddings of \(K\), and
\(YY^c=|\cdot|^{-1}\). Put
\[
           \Pi=\pi_0\otimes (Y|\cdot|^{1/2})\circ\det,
           \qquad \nu=V|_{G_H}\otimes\kappa_2.
\tag{G106N.12}
\]
Here the norm pullback of \(Y\) is understood at every
place of \(H\). The character \(Y|\cdot|^{1/2}\)
is unitary and conjugate self-dual. The same is true of
the resulting cuspidal parameter, since \(V\) is
conjugate self-dual on \(H\).

The normalization is the literal native \path{AttC}
normalization. At an unramified place of norm \(q\),
that definition multiplies the Satake roots in rank two
by \(q^{1/2}\). The scalar in (G106N.12) has value
\(\kappa_2(\operatorname{Fr}_{\rm geo})q^{-1/2}\).
Thus the resulting geometric-Frobenius roots are exactly
those of \(\nu\). At an embedding over the selected
embedding of \(K\), the infinitesimal exponents are
\(e-\tfrac12,-e-\tfrac12\); at its conjugate they
are their negatives in reverse order. They have the
required half-integral shift and are distinct. This
checks regular C-algebraicity, rather than asserting
that the untwisted unitary odd-weight parameter is
already C-algebraic. In native terminology
\(\nu\) satisfies \path{Af.IsAutCsd} with exactly
the retained attachment and cyclic base-change inputs.

The representation \(V\) factors through
\(G_\gamma\). At away inertia with \(\gamma=1\),
its actual residual-kernel killing condition makes the
inertia image finite prime-to-\(p\) with scalar residue
one; the scalar lemma makes it the identity.
Its polarization is symmetric with the required sign.
Indeed put
\(J_{\rm alt}=\left(\begin{smallmatrix}0&1\\-1&0\end{smallmatrix}\right)\)
and \(D=V(c)=\operatorname{diag}(-1,1)\).
For \(g\in G_H\), \(\det V(g)=1\) and
\(V(cgc^{-1})=DV(g)D\), so
\[
 V(g)^t(J_{\rm alt}D)V(cgc^{-1})=J_{\rm alt}D=J_2.
\tag{G106N.13}
\]
Thus it gives the symmetric rank-two native form, not
the alternating possibility. In this basis the
normalization difference at any \(g_0\) is already
diagonal, as in (G106N.9). The ordered local flag gives
the required ordered product relation. The continuous
integral matrices and local characters have their actual
strict residual coordinates from AC-G68, and therefore
give an actual point of the rational universal ring.
If its coefficient residue field is larger than \(k\),
use \(\mathcal O+\mathfrak m_E\subset\mathcal O_E\)
as the complete local test ring, as in AC-K29-point.
Composing this point with \(\Phi_2\) gives
\(x:\mathscr B_2\to\overline{\mathbb Q}_p\).
Its determinant and ordered local characters are those
of \(\nu\), with no replacement by a generic
pseudorepresentation.

We check all clauses of native \path{PT} at this point.
Continuity, the actual determinant realization, and
cuspidal automorphy were just proved. The auxiliary
place \(v_0\) was selected outside the bad places of
\(V\) and outside \(3p\), so \(\nu\) is unramified
there. At each coefficient-prime place, the two ordered
inertia characters of \(V\) are
\(\psi_p\epsilon^e\) and
\(\psi_p^{-1}\epsilon^{-e}\), up to their unramified
factors, with \(\psi_p^2=1\). The finite factors are
retained. Let \(b_\kappa(\tau)\in\{0,1\}\) be
the native local algebraic weight of \(\kappa_2\)
at an embedding \(\tau\) of the local field.
The two native local algebraic weights of \(\nu\) are
\[
                b_\kappa(\tau)-e,
                \qquad b_\kappa(\tau)+e.
\tag{G106N.14}
\]
They are strictly increasing. In the native
\path{IsWtPart} definition the finite characters belong
to its finite-inertia factor \path{ph}; they are not
discarded. The local reciprocity calculation in
\path{DiagonalCharactersRegularAlgebraicOfInertiaWeigh}
proves (G106N.14) by multiplication with that finite
factor, exactly as it proves the pure-cyclotomic case.
The possible inertia exponents of a quotient of two
ordered characters are \(0,2e,-2e\). None is one.
The cyclotomic character on inertia of any finite local
extension has infinite image, so none of the ordered
characters equals another times the native cyclotomic
character. This checks the last \path{PT} clause,
including the two equal-index cases. The numerical
range (G106N.1) in particular has \(e\ge2\).

Apply the retained \path{ModU} to the actual
\path{HidaSetting}, \path{Lev}, \path{Resid} and
native \path{LamObj} in rank two. Its literal statement
has no abelian-over-\(\mathbb Q\) hypothesis on the CM
host. The host has no \(\mu_p\), as proved in AC-G97,
and the chosen coefficient ring is large for it.
It supplies \(J_2\) with the complete native
\path{Modular} record. The just verified \path{PT}
conditions give \(J_2\subset\ker x\), hence
\(J_2\ne\mathscr B_2\). The same actual point,
with its rational normalized representation, is exactly
the source \path{IdealSpecializationPoint}, not merely
a point with the same semisimplified residual trace.

\proofstep{5. Every remaining field of the rank-two modular input.}
The native \path{Modular} record contains
\path{RingProps.fin}, \path{RingProps.tf},
\path{RingProps.unr0}, and \path{P1core}.
The native \path{LamEx} identifies its weight algebra
with \(\mathcal O[[X_1,\ldots,X_{2d}]]\), where
\(d=[F:\mathbb Q]\).
Thus \(\mathscr B_2/J_2\) is finite and torsionfree
over this domain and is nonzero. Every minimal prime
contracts to zero: otherwise a nonzero element of the
weight algebra would kill a nonzero element after
localizing at a minimal prime and then clearing a
denominator, contradicting torsionfreeness.
Integral finite dimension equality gives
\[
   \dim\mathscr B_2/Z=2d+1
           \quad\text{for every }Z\in\operatorname{Min}(J_2).
\tag{G106N.15}
\]
This is precisely the source \path{DimT} condition;
no flatness or reducedness of the whole quotient has
been added.

At the auxiliary decomposition group \(D_0\),
\(\gamma=1\), so both residual characters in
(G106N.6) equal \(\bar\kappa_2\). On its inertia
they are one. The actual auxiliary element has
\(\bar\epsilon\ne1\), by AC-G89 and its host
transport. These are all three fields of
\path{Resid.V0ok}. They imply \path{N4} for every
native point, by the complete matrix argument of
\path{EveryPointSatisfiesFrameConditionAlt}. To recall
why this does not require semisimplicity, if \(M\)
is the point's matrix of that auxiliary element, its
residual characteristic polynomial is \((X-a)^2\)
with \(a\ne0\). An intertwiner \(Y\) for the
cyclotomic twist satisfies
\(\epsilon M Y=Y M\). Cayley--Hamilton gives
\(\operatorname{charpol}_M(\epsilon M)Y=0\).
The matrix polynomial on the left is invertible,
because its reduction is a nonzero scalar plus a
nilpotent and \(\bar\epsilon\ne1\).
Thus \(Y=0\). This is the native \path{N4}
condition. The \path{unr0} field of \path{RingProps}
now gives the exact source \path{Unr0}, including
killing in the determinant algebra.


\proofstep{6. Literal native transfer and propagation.}
Surjectivity of \(\Phi_2\), native finite weight
support and \path{FinLamOf.lean} give
\(\operatorname{FinLam}(M_2)\).
One can also see the weight fiber directly. In
\(D=B_2/(M_2+\lambda_2)\), every rational local
character factors through the finite prime-to-\(p\)
quotient \(I_v/(I_v\cap\ker\gamma)\).
It therefore equals its residual character even over
the possibly nonreduced ring \(D\), by the scalar
lemma. The native characters (G106N.7) reduce to
their fixed residual values since the uniformizer is
zero in \(D\). The native \path{LamObj} universal
property consequently makes its map to \(D\) factor
through \(k\). As \(D\) is finite over that weight
algebra, it is finite over \(k\). Its only maximal
ideal is the image of \(\mathfrak m_{B_2}\), which
proves the radical assertion in (G106N.3).
The ideal \(M_2\) is proper because \(\Phi_2\) is
local and \(J_2\) is proper.

All hypotheses of the literal theorem \path{p1Of}
in \path{FirstPrimeConditionTransfer.lean} have now
been supplied. Its \path{pointEx} constructs the
native point from each one-dimensional local curve,
using the just constructed \path{TwLink}. Its
\path{niceOf} applies to the rational frame, where
\(\gamma\) is trivial on the local kernel of the
full cyclotomic character. Thus a nontrivial local
ratio on that kernel has no finite-image exception;
the changed-tuple cubic collision issue is not
present in this rank-two application. Together
with the verified \path{N4} and native
\path{P1core}, it gives exactly
\(\operatorname{P1'}(M_2)\) for all relevant
local primes.

The complete theorem \path{seedOf} in
\path{SeedOf.lean} requires only the actual
\path{Setup}, \path{IsPhi}, \path{Unr0},
\path{DimT} and \path{IdealSpecializationPoint}.
These have all been checked above. It supplies
\(\operatorname{SeedIn}(M_2)\) of dimension
\(2d+1\). Its proof passes through a minimal
native component below the actual seed, obtains
the generic irreducible representation from its
determinant, descends it across the exact
\path{TwLink} kernel, and uses \path{of_trace}
and \path{point_field} to recover the actual
normalized \(A_2,B_2\) maps. In rank two its
\path{of_trace} proof excludes the alternating
form by specialization to the symmetric
irreducible seed. Hence no odd-rank argument
has been substituted in this step. Native
universality gives its exact identity
\(\beta\Phi_2=q\), and therefore the full
component dimension, rather than a dimension
bound for an abstract invariant ring.

Apply the complete literal theorem
\path{components_good_rank_two} in
\path{ComponentsGoodRankTwo.lean} to the
\path{Ok}, \path{Big}, \path{FinLam},
\path{P1'} and \path{SeedIn} just obtained.
Its conclusion is \path{AllGood} for \(M_2\).
For rank two, \path{TI} is precisely
\path{OffTri}, since the additional rank-three
condition is vacuous. Thus \path{AllGood}
and \path{FinLam} are exactly (G106N.3),
including the quantifier over every off-triangular
prime. Finally apply AC-G106-even to this
particular \(M_2\). Its construction retains
the middle-character factor and proves the
changed-tuple finite weight fiber, so it gives
the asserted uniform ideal \(M_E\) without
setting the middle character equal to one in
the actual even ring.
\end{proof}

\paragraph{Source and geometric scope.}
All source paths in this paragraph have the pinned root
\path{Development/AutomorphicArithmetic/AutomorphicVarieties/}
at commit
\path{7cbf7a0c32fca1badf05931879a37c3f8a2b03b4}.
The native records are in
\path{WeilGroupCharacters/UniversalOrdinaryDeterminantRing.lean}
and \path{WeilGroupCharacters/ArithmeticPointClassicality.lean}.
The exact generic records and transfer theorems are in
\path{CurveDVRSpecialization/PropagationTheoremGL2.lean},
\path{SecondThirdPrimeConditionTransferAlt.lean},
\path{FinLamOf.lean}, \path{FirstPrimeConditionTransfer.lean},
\path{PointEx.lean}, \path{NiceOf.lean},
\path{SeedOf.lean} and \path{ComponentsGoodRankTwo.lean}
in that same curve-specialization directory.
The full applied containing proofs were read; the
source receipt distinguishes those reads from retained
upstream comparison, modularity, reciprocity and
automorphic base-change theorems. In particular this
section claims no independent Lean build or new proof
of \path{ModU}.

Geometrically, (G106N.2) is the actual morphism from the
normalized rational rank-two deformation chart to its
native ordinary determinant space. The single ideal
\(M_2\) is the pullback of native automorphic support;
the source propagation theorem extends its containment
from the proved seed component to all off-triangular
components. AC-G106-even then transports that same
support through the exact endpoint factor, with the
middle rank-one deformation still present. This uses
the already audited formal-deformation and finite-cover
dictionary; it introduces no new archival or motivic
vanishing assertion.



